\subsection{Publicación de inmuebles y ciclo de reservas}
\label{subsec:publicacion_reservas}

El módulo inmobiliario y de contratación rige el ciclo de vida mediante el cual los propietarios publican sus anuncios y los estudiantes formalizan solicitudes de alquiler vinculadas a un depósito o fianza.

\subsubsection{Formulario multi-paso y gestión multimedia en frontend}

Para maximizar la tasa de finalización del alta de una vivienda, el cliente web estructura el proceso a través del componente AccommodationFormPage.tsx mediante un asistente por etapas (*Multi-Step Wizard*):

\begin{itemize}
    \item \textbf{Etapa 1 -- Localización y geocodificación}: El anfitrión introduce los datos postales de la vivienda. La interfaz valida la dirección en tiempo real y muestra una previsualización de confirmación sobre el mapa interactivo.
    \item \textbf{Etapa 2 -- Características físicas y comodidades}: Se registran la superficie en metros cuadrados, el número de dormitorios, baños disponibles y el equipamiento del piso mediante selectores accesibles (*tags* de selección múltiple).
    \item \textbf{Etapa 3 -- Galería fotográfica y selección de portada}: El usuario puede subir múltiples fotografías en alta definición. Mediante la tabla asociativa listing_image_selection, el sistema permite ordenar las imágenes mediante arrastrar y soltar (*drag and drop*) y elegir cuál actuará como imagen principal de portada para el catálogo público.
    \item \textbf{Etapa 4 -- Condiciones contractuales y fianza}: Se fija el precio mensual del alquiler y el importe obligatorio de la fianza legal en garantía de conservación de la vivienda.
\end{itemize}

\subsubsection{Máquina de estados de reserva y control de concurrencia en backend}

La interacción contractual entre inquilino y anfitrión está gobernada por una máquina de estados estricta implementada en BookingRequestService:

\begin{enumerate}
    \item \textbf{Transiciones de estado formales}: Una solicitud de reserva nace en estado PENDING al ser enviada por el estudiante. El propietario dispone de un plazo temporal acotado para tomar una decisión: si la acepta, transiciona a ACCEPTED; si la declina, pasa a REJECTED; y si expira el plazo máximo sin respuesta, un proceso programado la marca automáticamente como EXPIRED.
    \item \textbf{Prevención de sobre-reserva mediante bloqueo optimista (\emph{Optimistic Locking})}: Para evitar que dos estudiantes reserven simultáneamente la misma habitación disponible en un entorno concurrido, la entidad Accommodation incorpora un campo de versión anotado con @Version. Si dos transacciones intentan confirmar la reserva al mismo instante, Hibernate detecta la colisión de versiones arrojando una excepción de concurrencia optimista (OptimisticLockException), abortando la segunda operación de forma segura y protegiendo la integridad del inventario.
    \item \textbf{Bloqueo y custodia de fianza}: Cuando una solicitud es aceptada, el sistema genera el registro financiero de la fianza en la tabla booking_requests, asociando un identificador de transacción y asegurando que las partes cuenten con un respaldo económico antes de formalizar la entrega de llaves.
\end{enumerate}

